Disaggregated systems separate compute, memory, and storage into independent network-attached services (e.g., Aurora, Snowflake, serverless platforms). Compute workers are instantiated for short durations, executing stateless functions without persistent local state. Storage and memory servers maintain persistent state accessed over network interconnects. This separation creates an impedance mismatch: state-dependent computation must query remote state over network links with latencies orders of magnitude higher than local memory access. Every validation query that requires consulting remote state introduces a network round-trip cost that dominates application latency.

The mismatch manifests in two distinct operations. The first is membership testing: determining whether an element exists in a remote set before executing read-modify-write operations. A compute worker querying a deterministic key-value store incurs a network round-trip for every validation. The second is cardinality estimation: computing the number of distinct elements in a partitioned dataset \cite{Harmouch2017CardinalityEA}. An exact \texttt{COUNT(DISTINCT)} query requires a shuffle operation that transmits all unique elements to a central reducer \cite{Zaharia2012ResilientDD}. Communication cost scales with dataset cardinality, rendering interactive analytics over petabyte-scale data infeasible.

Resolving this mismatch requires compact summaries that satisfy three conditions. The summaries must be queryable locally or transmitted with fixed cost independent of dataset size. They must synchronize with datasets that undergo insertion and deletion. They must support merge operations that are order-independent and produce results identical to processing the complete dataset when aggregating information across partitions.

Probabilistic data structures provide queryable summaries with sublinear space complexity. A Bloom filter occupies space proportional to the logarithm of the target false positive rate, independent of element size \cite{Bloom1970SpacetimeTI}. A HyperLogLog sketch estimates cardinality using fixed space regardless of dataset size \cite{Flajolet2007HyperLogLogTA}. These structures trade bounded error probability for space efficiency.

The architectural value of probabilistic data structures in disaggregated systems derives from two properties that address different subsets of the three conditions above. Dynamic mutability addresses state synchronization: structures exhibiting this property support element deletion in constant time without auxiliary counters, maintaining consistency with datasets that undergo explicit removal operations. Composability addresses distributed aggregation: structures exhibiting this property support merge operations that are commutative, associative, and idempotent, enabling distributed reduction without central coordination or data shuffling.

In \S\ref{sec:framework} we prove that these properties impose conflicting requirements on internal state representation. Mutability requires retaining sufficient information to identify and remove a specific element's contribution without affecting other elements. Composability requires information loss through commutative operations that discard element identities and enable order-independent aggregation.

This conflict explains why specific structures optimize for different distributed contexts. Cuckoo Filters satisfy dynamic mutability through explicit fingerprint storage in deterministically computable locations, enabling targeted deletion without introducing false negatives for remaining elements. HyperLogLog satisfies composability through register-wise maximum operations that preserve invariant properties regardless of aggregation order, enabling hierarchical reduction where partial results from multiple partitions combine into global aggregates.

\subsection{Organization}

Section~\ref{sec:membership} analyzes membership testing in distributed systems. We derive the requirement for dynamic mutability from state synchronization constraints and examine the Cuckoo Filter as a structure satisfying this property. We document deployment patterns enabled by mutability, including co-partitioned architectures, capacity scaling through consistent hashing, and concurrency control mechanisms.

Section~\ref{sec:cardinality} analyzes cardinality estimation in data-parallel systems. We derive the requirement for composability from distributed aggregation constraints and examine HyperLogLog as a structure satisfying this property. We document deployment patterns enabled by composability, including hierarchical reduction in MapReduce and stream processing systems, fault-tolerance protocols based on approximate recovery, and challenges arising from adversarial manipulation and privacy leakage in federated environments.

Section~\ref{sec:framework} formalizes the mutability-composability framework. We prove that idempotent merge operations and targeted deletion without false negatives are incompatible when reversible element representations are absent. We classify workloads by data lifecycle dynamics and query result formation, deriving decision criteria that map each workload class to the required property and corresponding architectural pattern.

Section~\ref{sec:conclusion} synthesizes the framework implications for disaggregated system design. We identify research directions at the intersection of the core properties with Byzantine resilience, differential privacy, and adaptive fault tolerance.
